\chapter{Deployment and Configuration}

\section{Hosting Architecture Overview}
The live production deployment operates across two decoupled, zero-cost cloud platforms:
\begin{itemize}[leftmargin=*]
    \item \textbf{Frontend}: Hosted globally on \textbf{Vercel Edge Network} (\url{https://bengaluru-traffic-ai.vercel.app/}).
    \item \textbf{Backend REST Service}: Hosted on \textbf{Render Web Services} (\url{https://bengalurutrafficai.onrender.com/}).
\end{itemize}

\section{Vercel SPA Rewrite Configuration}
React Router implements client-side dynamic routing across \texttt{/}, \texttt{/analytics}, and \texttt{/driver}. In traditional static web hosting, direct navigation or browser refresh on deep links triggers HTTP 404 Not Found errors because the static server looks for physical files at \texttt{/analytics/index.html}.

To solve this, \texttt{frontend/vercel.json} specifies automated server-side URL rewrites:
\begin{lstlisting}[language=json,caption={Vercel SPA Rewrite Rules (\texttt{frontend/vercel.json})}]
{
  "buildCommand": "npm install && npm run build",
  "outputDirectory": "dist",
  "framework": "vite",
  "rewrites": [
    {
      "source": "/(.*)",
      "destination": "/index.html"
    }
  ]
}
\end{lstlisting}
Every non-asset request is rewritten to \texttt{/index.html}, allowing the React Router bundle to parse the URL path and mount the corresponding view.

\section{Render Container Configuration}
Render manages the backend Python web service via the root blueprint specification:
\begin{lstlisting}[language=yaml,caption={Render Service Configuration (\texttt{render.yaml})}]
services:
  - type: web
    name: bengaluru-traffic-ai-backend
    runtime: python
    buildCommand: pip install -r requirements-demo.txt
    startCommand: python backend/app_demo.py
    envVars:
      - key: DB_NAME
        value: bengaluru_traffic
\end{lstlisting}
By configuring \texttt{requirements-demo.txt} and \texttt{backend/app\_demo.py}, the service builds and starts in under 45 seconds, requiring less than 85\,MB of RAM and operating well within Render's 512\,MB free limit.

\section{Environment Template and Secret Management}
Configuration settings are managed through environment variables without hardcoded credentials. Table~\ref{tab:env_vars} catalogs all supported variables documented in \texttt{.env.example}.

\begin{table}[h!]
\centering
\small
\begin{tabularx}{\textwidth}{l p{4.2cm} X}
\toprule
\textbf{Variable} & \textbf{Default Setting} & \textbf{Operational Function} \\
\midrule
\texttt{PORT} & \texttt{5000} & Network port for Flask HTTP server. \\
\texttt{FLASK\_ENV} & \texttt{development} & Server debug and reload mode flag. \\
\texttt{MONGO\_URI} & \texttt{mongodb://localhost:27017/} & MongoDB connection string (falls back to console if offline). \\
\texttt{DB\_NAME} & \texttt{bengaluru\_traffic} & Target database name for violation logs. \\
\texttt{VITE\_API\_URL} & \texttt{http://localhost:5000} & API base URL used by Vite frontend in production/development. \\
\texttt{JUNCTION\_NAME} & \texttt{Silk Board, Bengaluru} & Custom junction identity string displayed across UI. \\
\texttt{JUNCTION\_LAT} & \texttt{12.9176} & Reference latitude for Haversine distance calculations. \\
\texttt{JUNCTION\_LON} & \texttt{77.6238} & Reference longitude for Haversine distance calculations. \\
\bottomrule
\end{tabularx}
\caption{Environment Configuration Specification (\texttt{.env.example})}
\label{tab:env_vars}
\end{table}
